% AUTO-GENERATED by the update_record tool. DO NOT EDIT.
% verify_paper regenerates this file from record.json and the run outputs
% and fails on any difference, so manual edits always surface.
\noindent\textbf{Hypothesis 1.} SciGym 論文の Table 1 のうち Gemini-2.5-Pro、Gemini-2.5-Flash、GPT-4.1、GPT-4.1-mini の 4 行は、公開コードの Controller をそのまま用い、現在 API で提供されている同名モデル（Gemini 2 種は GA 版、GPT-4.1 2 種は論文と同じ 2025-04-14 スナップショット）で SciGym-small 137 件を最大 20 反復で解かせたとき、各モデルの平均 Simulation Trajectory Error（STE）が論文値から 0.05 以内に収まる形で再現される。~\cite[s1.p1, s1.p5, s1.p10]{duan-2025-measuring}
\begin{enumerate}
\item[\textbf{Claim 1}] 現行の Gemini-2.5-Pro（google/gemini-2.5-pro）で SciGym-small 137 件を最大 20 反復で解かせたときの平均 STE は、論文の 0.3212 を 0.05 より大きくは上回らない。~\cite[s1.p1, s1.p2, s1.p3, s1.p4, s1.p5, s1.p7, s1.p10]{duan-2025-measuring}, \cite[s2.p1]{h4duan-2025-scigym}
  \emph{Rationale:} Table 1 で最良とされた行がそのまま再現されることは、論文の主結果が特定のプレビュー版や実行環境に依存していないことの直接の証拠であり、h1 の Gemini-2.5-Pro 行を担う。
  \emph{Criterion:} \texttt{\detokenize{gemini-2.5-pro}}.\texttt{\detokenize{ste}} $-$ 0.3212 $\leq 0.05$.
  \emph{Prediction:} $[-0.05, 0.05]$ (論文 Table 1 の 1 回実行の平均値。137 件平均の標準誤差はおよそ 0.02 と見積もり、モデル版の差を含めても ±0.05 に収まると予測する).
  \emph{Observed:} 0.0935808.
  \emph{Verdict:} refuted.
\item[\textbf{Claim 2}] 現行の Gemini-2.5-Flash（google/gemini-2.5-flash）で SciGym-small 137 件を最大 20 反復で解かせたときの平均 STE は、論文の 0.4181 を 0.05 より大きくは上回らない。~\cite[s1.p1, s1.p2, s1.p3, s1.p4, s1.p5, s1.p6, s1.p10]{duan-2025-measuring}, \cite[s2.p1]{h4duan-2025-scigym}
  \emph{Rationale:} Gemini 系の mini 変種の行が再現されることで、h1 の Gemini-2.5-Flash 行を担い、c1 と合わせて Gemini 系の pro / mini の関係が論文どおりかを読めるようにする。
  \emph{Criterion:} \texttt{\detokenize{gemini-2.5-flash}}.\texttt{\detokenize{ste}} $-$ 0.4181 $\leq 0.05$.
  \emph{Prediction:} $[-0.05, 0.05]$ (論文 Table 1 の 1 回実行の平均値。137 件平均の標準誤差はおよそ 0.02 と見積もり、モデル版の差を含めても ±0.05 に収まると予測する).
  \emph{Observed:} 0.0284183.
  \emph{Verdict:} supported.
\item[\textbf{Claim 3}] GPT-4.1（openai/gpt-4.1、2025-04-14 スナップショット）で SciGym-small 137 件を最大 20 反復で解かせたときの平均 STE は、論文の 0.4611 を 0.05 より大きくは上回らない。~\cite[s1.p1, s1.p2, s1.p3, s1.p4, s1.p5, s1.p10]{duan-2025-measuring}, \cite[s2.p1]{h4duan-2025-scigym}
  \emph{Rationale:} 論文と同じ日付スナップショットが使える唯一の pro モデルであり、モデル版の差を含まない最も純粋な再現条件として h1 の GPT-4.1 行を担う。
  \emph{Criterion:} \texttt{\detokenize{gpt-4.1}}.\texttt{\detokenize{ste}} $-$ 0.4611 $\leq 0.05$.
  \emph{Prediction:} $[-0.05, 0.05]$ (論文 Table 1 の 1 回実行の平均値。同じスナップショットなので差は実行の揺らぎとシミュレータ差のみで、137 件平均の標準誤差およそ 0.02 の範囲に収まると予測する).
  \emph{Observed:} 0.0469407.
  \emph{Verdict:} supported.
\item[\textbf{Claim 4}] GPT-4.1-mini（openai/gpt-4.1-mini、2025-04-14 スナップショット）で SciGym-small 137 件を最大 20 反復で解かせたときの平均 STE は、論文の 0.6007 を 0.05 より大きくは上回らない。~\cite[s1.p1, s1.p2, s1.p3, s1.p4, s1.p5, s1.p6, s1.p10]{duan-2025-measuring}, \cite[s2.p1]{h4duan-2025-scigym}
  \emph{Rationale:} 同じスナップショットが使える mini モデルの行を担い、c3 と合わせて OpenAI 系の pro / mini の関係が論文どおりかを読めるようにする。
  \emph{Criterion:} \texttt{\detokenize{gpt-4.1-mini}}.\texttt{\detokenize{ste}} $-$ 0.6007 $\leq 0.05$.
  \emph{Prediction:} $[-0.05, 0.05]$ (論文 Table 1 の 1 回実行の平均値。同じスナップショットなので差は実行の揺らぎとシミュレータ差のみで、137 件平均の標準誤差およそ 0.02 の範囲に収まると予測する).
  \emph{Observed:} 0.00101829.
  \emph{Verdict:} supported.
\end{enumerate}
\noindent\emph{Assumed, so that the claims together imply Hypothesis 1:}
\begin{itemize}
\item 同名モデルの現行版（Gemini-2.5-Pro / Flash の GA 版、GPT-4.1 / 4.1-mini の 2025-04-14 スナップショット）が、論文で使われた版と同等の能力を持つこと（c1〜c4。Gemini はプレビュー版が提供終了のため GA 版で代替する）
\item aarch64 向けの libroadrunner 2.7.0 と antimony 2.14.0 が、論文環境の libroadrunner 2.8.0 と同じ軌道と評価値を与えること（c1〜c4）
\item temperature 1.0 での 1 回実行の平均 STE（137 件平均）の実行間変動が 0.05 より十分小さいこと（c1〜c4）
\item Vercel AI Gateway の OpenAI 互換エンドポイント経由の各モデルが、各社の公式 API と同じ応答分布を持つこと（c1〜c4）
\item 平均 STE が Table 1 の再現性を代表すること。RMS（modifier あり / なし）は同じ実行から表として報告するが、判定基準には用いない（c1〜c4）
\end{itemize}
